\subsection{B.1. Minimal Cohen factorization and the pro-Tor comparison}

\begin{theorem}[Known, with the manuscript form checked]\label{SA23-B1-1}
Let $(A,\mathfrak m,k)\to(B,\mathfrak n,K)$ be a local homomorphism
of Noetherian local rings, and suppose $B$ is complete.  There is a
factorization $A\to C\twoheadrightarrow B$ by local homomorphisms,
where $C$ is complete Noetherian local, $A\to C$ is flat, $C_0=C/\mathfrak mC$
is regular, and
\[
 \operatorname{gr}(C_0)\longrightarrow\operatorname{gr}(B/\mathfrak mB)
\]
is an isomorphism in degrees zero and one.  Associated graded rings
here use the maximal-ideal filtrations of the two closed fibers.
\end{theorem}
\begin{proof}[Derivation from the opened theorem]
Nasseh--Sather-Wagstaff, \emph{Cohen factorizations: Weak functoriality
and applications}, arXiv:1205.2119v3, Definitions 2.3 and 2.5 and
Fact 2.6(a), state existence of a minimal Cohen factorization.
They identify the original result as Avramov--Foxby--Herzog,
\emph{Structure of local homomorphisms}, Theorem (1.1) and
Proposition (1.5).  We opened the former full text; the latter original
is recorded separately as not opened.

In these definitions a Cohen factorization has precisely the
completeness, surjectivity, flatness, and regular-closed-fiber properties
above.  Its minimality says
$\dim C-\dim A=\operatorname{edim}(B/\mathfrak mB)$.
For a flat local map the dimension formula gives
$\dim C-\dim A=\dim C_0$, which equals
$\operatorname{edim}C_0$ because $C_0$ is regular.  The surjection
$C_0\twoheadrightarrow B/\mathfrak mB$ identifies residue fields with
$K$ and induces a surjective map of finite-dimensional $K$-vector
spaces on cotangent spaces.  Their dimensions agree by minimality,
so the map in degree one is bijective.  Degree zero is the residue
field isomorphism.  Conversely those two graded isomorphisms force
the same embedding-dimension equality and hence minimality.  This
establishes the precise identification with Deligne's scan 061 lemma,
without adding a separability assumption on $K/k$.
\end{proof}

\begin{lemma}[Proved: finite-complex Artin--Rees comparison]\label{SA23-B1-2}
Let $R$ be Noetherian, $J\subset R$ an ideal, and $L_\bullet$ a
complex whose terms are finite $R$-modules.  For each fixed degree
$p$, the canonical map of pro-$R$-modules
\[
 \{H_p(L)/J^rH_p(L)\}_{r\ge1}
 \longrightarrow \{H_p(L/J^rL)\}_{r\ge1}
\]
is an isomorphism.
\end{lemma}
\begin{proof}
Write $Z_p=\ker(d:L_p\to L_{p-1})$ and
$D_p=\operatorname{im}(L_{p+1}\to L_p)$.  Artin--Rees supplies a
constant $c$ such that
$Z_p\cap J^{r+c}L_p\subset J^rZ_p$ and
$D_{p-1}\cap J^{r+c}L_{p-1}\subset J^rD_{p-1}$ for all $r$.
If a class represented by $z\in Z_p$ becomes a boundary modulo
$J^{r+c}L$, write $z=b+w$ with $b\in D_p$ and
$w\in Z_p\cap J^{r+c}L_p$.  Its class in $H_p(L)$ belongs to
$J^rH_p(L)$.  Thus the kernels of the displayed maps form a zero
pro-object.

If $z\in L_p$ is a cycle modulo $J^{r+c}L$, then
$dz\in D_{p-1}\cap J^{r+c}L_{p-1}\subset J^rD_{p-1}$.
Choose $w\in J^rL_p$ with $dw=dz$.  The element $z-w$ is a genuine
cycle and has the same class as $z$ modulo $J^rL$.  Thus after the
transition from index $r+c$ to $r$, every class comes from $H_p(L)$;
the cokernels form a zero pro-object as well.  A map of inverse
systems is an isomorphism in the abelian category of pro-modules
when its kernel and cokernel are zero pro-objects.  The bounds above
prove that criterion directly in this case.
\end{proof}

\begin{theorem}[Proved: Deligne's pro-Tor assertion]\label{SA23-B1-3}
Let $A\to B$ be a local homomorphism of Noetherian local rings,
let $\mathfrak a$ be any ideal of $A$, let $N$ be a finite $A$-module,
and let $M$ be a finite $B$-module.  For every $p\ge0$ the canonical
map of pro-$B$-modules
\[
 \left\{\operatorname{Tor}^A_p(N,M)\big/
      \mathfrak a^r\operatorname{Tor}^A_p(N,M)\right\}_r
 \longrightarrow
 \left\{\operatorname{Tor}^A_p(N/\mathfrak a^rN,M)\right\}_r
\]
is an isomorphism.  In particular the inverse limit on the right is
the $\mathfrak a$-adic completion of $\operatorname{Tor}^A_p(N,M)$.
\end{theorem}
\begin{proof}
First suppose $B$ is complete and choose $A\to C\twoheadrightarrow B$
as in Theorem~\ref{SA23-B1-1}.  The module $M$ is finite over $C$.
Take a free $C$-resolution $F_\bullet\to M$ with finite free terms;
Noetherianity constructs it successively.  Since $C$ is $A$-flat,
it is also an $A$-flat resolution.  Thus the finite-$C$-module complex
$L=F\otimes_A N$ has homology
$H_p(L)=\operatorname{Tor}^A_p(N,M)$, and
\[
 L/\mathfrak a^rL=F\otimes_A(N/\mathfrak a^rN)
\]
has homology $\operatorname{Tor}^A_p(N/\mathfrak a^rN,M)$.
Apply Lemma~\ref{SA23-B1-2} over $C$ with $J=\mathfrak aC$.
The comparison is the canonical one induced by $N\to N/\mathfrak a^rN$.

For arbitrary $B$, the map $B\to\widehat B$ is faithfully flat and
$\widehat M=M\otimes_B\widehat B$ is finite over $\widehat B$.
Using an $A$-free resolution of $N$ with finite free terms shows that
Tor, its quotients, and the comparison maps commute with this flat
base change.  The completed case makes the kernel and cokernel
systems pro-zero after tensoring with $\widehat B$.  For each index
it supplies a later index at which the transition map is zero after
base change; faithful flatness makes that same transition map zero
before base change.  The original kernel and cokernel systems are
therefore pro-zero.  Taking inverse limits of the pro-isomorphism
gives the asserted completion identity.  The ideal need not be the
maximal ideal; this matches the explicit ``un id\'eal de $A$'' on
scan 061.
\end{proof}
